\documentclass[11pt]{article}
\usepackage{amsmath, amssymb, amsfonts}
\usepackage{geometry}
\usepackage{booktabs}
\geometry{margin=1in}

\title{\textbf{CSE400: Fundamentals of Probability in Computing}\\
Lecture 20: Linear Transformations and Expectations of Multiple RVs}
\author{Instructor: Dhaval Patel, PhD}
\date{}

\begin{document}

\maketitle

\section{Statistical Representation of Random Vectors}

\subsection{Definition}
For multiple random variables $X_1, X_2, \dots, X_N$, we define a random vector:
\begin{equation}
\mathbf{X} =
\begin{bmatrix}
X_1 \\
X_2 \\
\vdots \\
X_N
\end{bmatrix}
\end{equation}

\textbf{Purpose:} Compact representation for handling multiple random variables.

\subsection{Mean Vector}
\begin{equation}
\boldsymbol{\mu}_X = \mathbb{E}[\mathbf{X}]
=
\begin{bmatrix}
\mathbb{E}[X_1] \\
\mathbb{E}[X_2] \\
\vdots \\
\mathbb{E}[X_N]
\end{bmatrix}
\end{equation}

\textbf{Logic:} Represents average behavior of the system.

\subsection{Correlation Matrix}
\begin{equation}
\mathbf{R}_{XX} = \mathbb{E}[\mathbf{X}\mathbf{X}^T]
\end{equation}

\begin{equation}
R_{XX}(i,j) = \mathbb{E}[X_i X_j]
\end{equation}

\textbf{Property:} Symmetric $N \times N$ matrix.

\subsection{Covariance Matrix}
\begin{equation}
\mathbf{C}_{XX} = \mathbb{E}[(\mathbf{X} - \boldsymbol{\mu}_X)(\mathbf{X} - \boldsymbol{\mu}_X)^T]
\end{equation}

\textbf{Structure:}
\begin{equation}
\mathbf{C}_{XX} =
\begin{bmatrix}
\mathrm{Var}(X_1) & \mathrm{Cov}(X_1,X_2) & \cdots \\
\mathrm{Cov}(X_2,X_1) & \mathrm{Var}(X_2) & \cdots \\
\vdots & \vdots & \ddots
\end{bmatrix}
\end{equation}

\textbf{Property:} Symmetric since $\mathrm{Cov}(X_i,X_j)=\mathrm{Cov}(X_j,X_i)$.

\section{Linear Transformations of Random Vectors}

\subsection{Model}
\begin{equation}
\mathbf{Y} = \mathbf{A}\mathbf{X} + \mathbf{b}
\end{equation}

\begin{itemize}
    \item $\mathbf{A}$: transformation matrix
    \item $\mathbf{b}$: constant shift vector
\end{itemize}

\subsection{Mapping Properties}

\begin{center}
\begin{tabular}{ll}
\toprule
Property & Result \\
\midrule
Mean & $\boldsymbol{\mu}_Y = \mathbf{A}\boldsymbol{\mu}_X + \mathbf{b}$ \\
Covariance & $\mathbf{C}_{YY} = \mathbf{A}\mathbf{C}_{XX}\mathbf{A}^T$ \\
Correlation & $\mathbf{R}_{YY} = \mathbf{A}\mathbf{R}_{XX}\mathbf{A}^T + \mathbf{A}\boldsymbol{\mu}_X \mathbf{b}^T + \mathbf{b}\boldsymbol{\mu}_X^T \mathbf{A}^T + \mathbf{b}\mathbf{b}^T$ \\
\bottomrule
\end{tabular}
\end{center}

\subsection{Derivations}

\textbf{Mean:}
\begin{equation}
\mathbb{E}[\mathbf{Y}] = \mathbb{E}[\mathbf{A}\mathbf{X} + \mathbf{b}] = \mathbf{A}\boldsymbol{\mu}_X + \mathbf{b}
\end{equation}

\textbf{Covariance:}
\begin{equation}
\mathbf{Y} - \boldsymbol{\mu}_Y = \mathbf{A}(\mathbf{X} - \boldsymbol{\mu}_X)
\end{equation}

\begin{equation}
\mathbf{C}_{YY} = \mathbb{E}[\mathbf{A}(\mathbf{X} - \boldsymbol{\mu}_X)(\mathbf{X} - \boldsymbol{\mu}_X)^T \mathbf{A}^T]
= \mathbf{A}\mathbf{C}_{XX}\mathbf{A}^T
\end{equation}

\textbf{Note:} Shift $\mathbf{b}$ does not affect covariance.

\section{Multivariate Gaussian Distribution}

\subsection{1D Case}
\begin{equation}
f_X(x) = \frac{1}{\sqrt{2\pi\sigma^2}} \exp\left(-\frac{(x-\mu)^2}{2\sigma^2}\right)
\end{equation}

\subsection{N-Dimensional Case}
\begin{equation}
f_{\mathbf{X}}(\mathbf{x}) =
\frac{1}{(2\pi)^{N/2} |\mathbf{C}_{XX}|^{1/2}}
\exp\left(-\frac{1}{2}(\mathbf{x}-\boldsymbol{\mu}_X)^T \mathbf{C}_{XX}^{-1} (\mathbf{x}-\boldsymbol{\mu}_X)\right)
\end{equation}

\subsection{Uncorrelated Gaussian Variables}

\begin{itemize}
    \item $\mathrm{Cov}(X_i,X_j)=0$ for $i \neq j$
    \item $\mathbf{C}_{XX}$ is diagonal
\end{itemize}

\textbf{Factorization:}
\begin{equation}
f_{\mathbf{X}}(\mathbf{x}) = \prod_{n=1}^{N}
\frac{1}{\sqrt{2\pi\sigma_n^2}}
\exp\left(-\frac{(x_n - \mu_n)^2}{2\sigma_n^2}\right)
\end{equation}

\section{Worked Example}

Given:
\begin{equation}
\boldsymbol{\mu}_X =
\begin{bmatrix}
2 \\ -1 \\ 1
\end{bmatrix},
\quad
\mathbf{A} =
\begin{bmatrix}
1 & 0 & -1 \\
0 & -1 & 1 \\
-1 & 1 & 0
\end{bmatrix}
\end{equation}

\begin{equation}
\mathbf{R}_{XX} =
\begin{bmatrix}
13 & 2 & 3 \\
2 & 10 & 3 \\
3 & 3 & 10
\end{bmatrix}
\end{equation}

\subsection{Mean}
\begin{equation}
\boldsymbol{\mu}_Y = \mathbf{A}\boldsymbol{\mu}_X =
\begin{bmatrix}
1 \\ 2 \\ -3
\end{bmatrix}
\end{equation}

\subsection{Correlation}
\begin{equation}
\mathbf{R}_{YY} =
\begin{bmatrix}
17 & -6 & -11 \\
-6 & 14 & -8 \\
-11 & -8 & 19
\end{bmatrix}
\end{equation}

\subsection{Covariance}
\begin{equation}
\boldsymbol{\mu}_Y \boldsymbol{\mu}_Y^T =
\begin{bmatrix}
1 & 2 & -3 \\
2 & 4 & -6 \\
-3 & -6 & 9
\end{bmatrix}
\end{equation}

\begin{equation}
\mathbf{C}_{YY} =
\begin{bmatrix}
16 & -8 & -8 \\
-8 & 10 & -2 \\
-8 & -2 & 10
\end{bmatrix}
\end{equation}

\section{Case Study: Smart City Traffic (Ahmedabad)}

\subsection{Variables}
\begin{itemize}
    \item $X_1$: Traffic flow
    \item $X_2$: Speed
    \item $X_3$: Rainfall
    \item $X_4$: Temperature
    \item $X_5$: AQI
\end{itemize}

\subsection{Interpretation}
\begin{itemize}
    \item $\mathrm{Cov}(X_1,X_3) > 0$
    \item $\mathrm{Cov}(X_2,X_3) < 0$
    \item $\mathrm{Cov}(X_1,X_5) > 0$
\end{itemize}

\subsection{Congestion Index}
\begin{equation}
Y = 0.6X_1 - 0.4X_2 + 0.8X_3
\end{equation}

\textbf{Note:} This is a linear transformation ($Y = \mathbf{a}^T \mathbf{X}$).

If $\mathbf{X}$ is Gaussian, then $Y$ is also Gaussian.

\section{Homework}

\begin{itemize}
    \item Solve: $\boldsymbol{\mu}_X \boldsymbol{\mu}_X^T = \mathbf{R}_{XX} - \mathbf{C}_{XX}$
    \item Prove: $\mathrm{Var}(X+Y) = \mathrm{Var}(X) + \mathrm{Var}(Y) + 2\mathrm{Cov}(X,Y)$
    \item Derive 2D Gaussian PDF using determinant and inverse
\end{itemize}

\end{document}